% ========== 前言 ==========
\chapter*{前言}
\addcontentsline{toc}{chapter}{前言：从文生图到镜头控制}

当下以 Stable Diffusion 为代表的文生图扩散模型已经能够根据文本提示词生成质量极高的图像，但在镜头角度、空间结构、物体轮廓与多视角一致性等方面的控制能力依然有限：文本对镜头的控制本质上是不稳定的，例如输入"俯拍""近景""镜头向左移动"，模型往往只能生成一个"看起来差不多"的画面，却并未真正理解相机的位置、角度与视场等几何信息。换言之，普通 prompt 更多控制的是语义与风格，而非严格的空间结构。

本技术报告所记录的工作，目标不在于单张图片的随机生成，而在于：给定一张参考图与相机相对旋转的三维向量提示词，让扩散模型生成同一物体在目标视角下、且尽量保持结构、身份、材质与光照一致性的图像。为达成这一目标，我们将抽象的相机控制转化为模型可理解的结构控制——通过深度图与 Canny 轮廓图这两类中间条件图，再借助 ControlNet 的结构注入能力与 LoRA 的轻量化微调，使扩散模型在生成时受到明确的几何约束。

整套系统并非已经完全实现了严格的六自由度（6DoF）相机控制，更准确的定位是一个"基于深度图、Canny 轮廓图、LoRA 与 ControlNet 的镜头控制原型系统"。它验证了一个核心思路：通过中间结构条件，可以让扩散模型更稳定地服从镜头控制，从而为后续的多视角一致生成与连续镜头生成打下基础。

本报告的章节安排如下：第一章为绪论，交代问题背景、目标与贡献；第二章梳理相关工作；第三章详述整体方法设计；第四章介绍基于 Blender 的多视角数据集构建管线；第五章给出训练、收敛验证、推理应用与 ComfyUI 部署的实现细节；第六章对结果与局限进行分析讨论；第七章总结全文并展望未来工作。
